\section{Introduction}
\label{sec:intro}

An attacker who has just landed inside a honeypot issues a single
request: \emph{show me the database credentials}. What the honeypot does
next divides two decades of deception engineering from the problem this
paper addresses. A recorder cannot answer at all; it writes the request
to disk and waits. A scripted honeypot answers from a file, the same file
for every attacker, on every visit, forever. A generative honeypot
drafts an answer on the spot -- and whatever that draft contains, the
attacker will believe it, act on it, and carry it into a timeline that a
security operations center will later treat as evidence. If the draft
invents a credential, that credential exists nowhere except in the
attacker's hands, and nothing downstream can tell it apart from a secret
that leaked. If the draft echoes a real one, the honeypot has just
converted a trap into a breach. The moment a honeypot begins to answer,
it has taken on all three jobs at once: stay consistent with everything
it has said so far, leak nothing real, and remain available when the
model behind it is not. A recent survey of production deception frameworks
identifies a gap in the surveyed systems: responses are not grounded in
real-time threat intelligence and behavior is not differentiated across
attacker classes~\cite{lopez2024}. We treat that observation as scoped
related-work evidence, not a claim about every deployed honeypot.

The cost of that silence is measured elsewhere. Mean time to identify a
breach runs in
months~\cite{ibm2024}, while adversary breakout time has compressed to
minutes~\cite{crowdstrike2026}. A honeypot sits directly in the path of the
attacker, yet converts almost none of that proximity into understanding: it
cannot tell a credential-stuffing bot from a hands-on-keyboard operator,
and it produces none of the
tactics-and-techniques intelligence that a security operations center would
need to act on the encounter.

Large language models (LLMs) promise to close this gap. Prompted with a
session transcript and a library of threat intelligence, a modern LLM can
impersonate a vulnerable Linux host, improvise plausible file systems, and
steer the conversation toward the attacker's next objective. Recent systems
have begun to exploit this, deploying LLM agents as honeypots that chat an
attacker through a fictitious machine~\cite{reworr2024}. The promise,
however, arrives with a problem the deception literature has largely left
unexamined: \emph{the model can be wrong in ways a script cannot}. A
honeypot's product is a lie, and the lie must satisfy three properties at
once. It must be \emph{consistent}: the kernel version printed on turn one
cannot contradict the one an attacker reads on turn five, and the database
banner cannot disagree with the port the same host just advertised. It
must be \emph{safe to serve}: a model that complies too well will invent
credential material with no way to distinguish what it fabricated from
what it leaked, poisoning every downstream analytical product. And it
must be \emph{available}: a generator is a rented service, and ours was.
Operational traces in this study show catalog drift, quota exhaustion,
gateway outages, and safety-alignment refusals. Unverified generation can
forfeit all three properties in the same session.

A natural question arises: why not fix this inside the model? Careful
prompting, refusal training, and constrained decoding all shape
probabilities, but none can certify a particular output. Alignment
training produces models whose refusal behavior shifts with phrasing and
pressure~\cite{ganguli2022,bai2022}, which is precisely the failure our
quota cliff exposed; and a deceived honeypot is itself an LLM application
with attacker-controlled context, so injection techniques that hijack
system prompts~\cite{greshake2023} apply with the attacker already in the
room. A probability, however small, is not a promise a security tool
can serve on. The assurance has to live outside the model, in the
pipeline around it.

The classical system-side fixes fail differently. Pre-scripting every
answer regresses to the finite automaton of low-interaction deception,
discarding the adaptivity that motivated generation in the first place;
real-time human review of every response cannot keep pace with an
attacker who reads faster than an analyst can vet. What survives is an
asymmetry between the two halves of the problem. Generating a plausible
answer to an arbitrary question is open-ended: its output space cannot
be enumerated, so no amount of testing certifies the next draw.
Checking a draft against a small, closed set of claims is the opposite:
a bounded computation whose verdict is final. A system can exploit this
by splitting the roles, letting the generator propose and a verifier
dispose against a fiction the system itself fully controls.

This paper takes the position that the right question is not
\emph{can an LLM impersonate a host} but \emph{how much intelligence can
a lie safely spend, and where must the spending stop}. Deception systems
build intelligence in rungs -- templates, classifiers, retrieval, free
generation -- and each rung has a trust price that practice pays
implicitly. We make the pricing explicit. Our answer is a verified
deception loop in which
each mechanism exists to buy one of the three properties, placed on a
single axis of certifiability (\S\ref{sec:ladder}). A fiction graph
binds every host fact to one consistent set of triples, and a verifier
checks every draft against that graph before anything reaches the wire:
consistency, enforced at serving time, not asked for at prompt time.
Baited content is drawn only from canary-registered secrets, so that any
disclosure is a measured event rather than a leak: safety by
construction. A deterministic fallback serves human-vetted transcripts
whenever generation or verification fails, so the system degrades to the
last known good answer rather than an unverified improvisation:
availability under provider failure. Beneath these sit a behavioral
profiler and a hybrid retrieval engine that ground a per-attacker
deception plan in curated CTI, giving the lie its targeting, and a
human-in-the-loop gate holds back any content class marked as requiring
approval. The loop turns generation from an act of trust into a measured
pipeline with an audit trail, and the evaluation reads the whole
apparatus as an escalation study: each rung of intelligence is measured
for what it adds over the rung beneath it, including the rungs that fold
back down when a certifiable mechanism reproduces their effect.

We evaluate whether each element of this loop contributes measurably,
using
artifact-backed benchmarks that we release alongside the paper. Four
questions structure the study:
\begin{itemize}
\item \textbf{RQ1:} Does retrieval grounding measurably change CTI analysis,
and separately what is served? A paired, per-scenario archived ablation across
90 matched scenarios reports that grounding lifts analysis-channel technique detection by
$+0.317$ (95\% CI $[+0.162, +0.480]$) and actor attribution by $+0.140$
($[+0.083, +0.197]$), while leaving persona realism untouched, a null result
the architecture predicts. The same-model served-text comparison and
echo-guarded audit are reported separately; we do not interpret the raw served
score as a clean serving-quality causal result.
\item \textbf{RQ2:} Can an LLM-driven honeypot remain internally consistent
and safe under sustained adversarial interaction? Across 45 targeted
re-queries the fiction graph serves zero contradictions, and across a
218-query live campaign the verification layer records no unverified
response served. Separate induced-draft tests measure the verifier rather
than only the clean workload: the 105-case deterministic suite is perfect,
while three authored model families expose bounded precision/recall gaps.
\item \textbf{RQ3:} Is the resulting deception strong enough to hold a
skilled observer under an internal diagnostic? A three-evaluator consensus study, each evaluator an
independently prompted LLM reviewer persona distinct from the generator,
scores the full session transcripts at $3.86/5$ (exact agreement
$0.750$, a within-model consistency measure -- all three personas share
one judge model), with all eight scenarios at or above the rubric-midpoint
acceptance bar and
a known judge bias corrected by canary-registered secrets.
\item \textbf{RQ4:} Does each additional rung of intelligence justify its
trust price? The escalation analysis adds each mechanism
to the same
loop and prices its marginal effect: on one 400-session held-out split the
profiler scores $0.8875$ as a whole-session classifier and $0.955$ when the
trajectory replay exposes each session's full command history---the same
corpus under two measurement procedures, not an intervention delta;
grounding contributes the archived analysis-channel $+0.154$
comparison above; and at the top rung, kind-aware adversarial classification
corrects 12/13 recorded Cowrie-log assessments and raises attacker bait exposure by 125\% in a
five-pair closed-loop replay, yet adds no measurable dwell time and is
measurably subsumed by a deterministic decision floor at the benchmark's
resolution -- and a second intelligent rung (LLM-fusion retrieval) is
priced against deterministic BM25 and shows fold-down pressure on three
of four metrics -- a directional fold-down signal that turns the ladder
from a build history toward a design rule whose generality this
benchmark cannot yet certify.
\end{itemize}

\noindent Our contributions are:
\begin{itemize}
\item \textbf{A verified deception architecture} (\S\ref{sec:design}) that
couples RAG-grounded generation to a fiction graph and a serving-time
verification gate, with deterministic fallback, so that every served
response is either verified against the deception state or produced from
vetted static content.
\item \textbf{An archived analysis-channel grounding comparison}
(\S\ref{sec:evaluation}): paired-difference bootstrap intervals over
identical 90-row scenario sets report $+0.154$ overall improvement
($[+0.092, +0.221]$), with 35 TTP detections spanning 21 distinct techniques
versus zero without grounding. The archived artifact predates canonical
scenario IDs; the new rerun contract records them, and no clean served-text
causal gain is claimed: our own scorer audit yields echo-corrected
both-channel lower bounds of 0.479 technique and 0.243 actor for the archived
campaign.
\item \textbf{A consistency and safety benchmark} for LLM honeypots
(\S\ref{sec:evaluation}): contradiction-free serving under adversarial
re-query, next-technique steering accuracy (top-3 $0.9$), a
completeness-verifying supported-claim gate measured at precision and
recall $1.0$ over a 123-case suite whose declared surface is eight
commands, plus a 105-case multi-check
deterministic suite and a cross-family authored-draft stress test, a
host-keyed cross-session fiction layer that
closes the between-sessions gap at zero cost on benign traffic, and an
ablation battery in which removing each
enforcement mechanism (A1--A4) costs the system the property it was
added to provide, while the A5 check validates safe-path equivalence.
\item \textbf{Run-to-run stability evidence} rarely reported for LLM
systems: two archived 218-query reruns reproduce the composite to
$+0.005$ with $0.991$ per-row pass agreement, the deterministic
strata reproduce exactly, the comparison survives unchanged when
recomputed on the judge-audited rows, and per-serve provenance labels
attribute 216 of 218 responses to verified generation with no fail-safe
substitution.
\item \textbf{A release package} on commodity hardware
(2~vCPU, ${\sim}4$~GB RAM), in which every headline number resolves to a
persisted result file at a repository-root path, and three scripts
regenerate the derived intervals and the figure values rather than merely
storing them. The package is published on artifact acceptance; this anonymous
submission carries no locator, and provider-backed reruns are labelled
unavailable.
\item \textbf{An escalation methodology for AI in deception}
(\S\ref{sec:ladder}, \S\ref{sec:jeveval}): mechanisms placed on a
certifiability axis, each rung priced by measured marginal effect, and
the descent direction exercised on live machinery -- when a certifiable
mechanism reproduces an intelligent one's benefit, the intelligent one
comes out. The same axis also admits the converse result: a frozen
candidate mechanism (the Evidence Judge/filter) was tested on a
prospective 40-case cohort that exercised the actual HTTP and SSH
serving path and found to add no measurable value beyond the baseline
gate on that configuration, so it was retained behind its
disabled-by-default boundary rather than claimed.
\item \textbf{An audited measurement pipeline} (\S\ref{sec:evaluation}):
the automated judge itself is audited offline against the transcripts it
scored, a full-corpus census locates the stratum where the fold-down
comparison is informative, and the audit-plus-stability chain reseeds
automatically whenever a new campaign run lands.
\end{itemize}

We are careful about what this paper does not claim. The deployment is
tailnet-only, so it observes no ambient internet background radiation; the
evaluation judge shares a provider with, but is a different model from, the
generator; the quality score is not human validation; and the system's
availability on free-tier inference is a constraint any practitioner should
weigh. The served numbers come from the released default path on commodity
hardware; the full-system study that exercises the actual HTTP and SSH
serving path was run in deterministic provider-unavailable mode only, and
no provider-backed inference or external deployment is claimed from it.
The Evidence Judge/filter and the conditional Laya scorer are separate
experiment identities --- a frozen stage tested on a prospective cohort
and a separately-versioned feasibility study, respectively --- and neither
is folded into the served loop's measured contributions; their records
stand where they stand. The threat model excludes host compromise, provider
compromise, and operator collusion, and the closed command surface does not
establish open-world robustness. These limits are quantified where they
arise (\S\ref{sec:threat-model}, \S\ref{sec:conclusion}).

The remainder of the paper is organized as follows.
Section~\ref{sec:threat-model} fixes the threat model and claim boundary;
Section~\ref{sec:design} presents the verified deception loop.
Section~\ref{sec:related} positions it against generative honeypots,
runtime enforcement, retrieval for threat intelligence, and the attacks on
that retrieval layer.
Section~\ref{sec:evaluation} reports the benchmark battery, the ablations,
the consensus study, and the stability runs; Appendix~\ref{app:artifacts}
indexes the released artifacts behind every number.
Section~\ref{sec:conclusion} discusses limitations, deployment guidance,
and the conclusion.
